% --- Archon physics formalization source begin ---
% archon:physics
% archon:covers PhyXMiniProblems/problem_phyx_mini_0711.lean
% archon:source-report reports/phyx_mini/problem_phyx_mini_0711.source.json

\chapter{Physics problem phyx\_mini\_0711}
\label{ch:PhyXMiniProblems_problem_phyx_mini_0711}

\paragraph{Problem source.}
\#\# Physical scenario

Two equal masses are at the ends of a massless \textbackslash{}(50\textbackslash{},\textbackslash{}mathrm\{cm\}\textbackslash{})-long rod. The rod spins at \textbackslash{}(2.0\textbackslash{},\textbackslash{}mathrm\{rev/s\}\textbackslash{}) about an axis through its midpoint. Suddenly, a compressed gas expands the rod out to a length of \textbackslash{}(160\textbackslash{},\textbackslash{}mathrm\{cm\}\textbackslash{}).

\#\# Question

What is the rotation frequency after the expansion?

\#\# Answer choices

- A: A: 0.80rev/s

- B: B: 0.50rev/s

- C: C: 0.20rev/s

- D: D: 1.10rev/s

\#\# Figure caption (auxiliary; use the image as primary evidence)

The image consists of two diagrams labeled "Before" and "After," illustrating a rotational motion concept.

**Before:**

- A rod with two attached spheres labeled 1 and 2 is rotating with an initial angular velocity (\textbackslash{}( \textbackslash{}omega\_i = 2 \textbackslash{}, \textbackslash{}text\{rev/s\} \textbackslash{})).
- The rod length (\textbackslash{}( l\_i \textbackslash{})) is 50 cm.
- Arrows represent linear momenta (\textbackslash{}( \textbackslash{}vec\{p\}\_\{1i\} \textbackslash{}) and \textbackslash{}( \textbackslash{}vec\{p\}\_\{2i\} \textbackslash{})) for spheres 1 and 2, pointing outward and downward/leftward respectively.
- Total angular momentum (\textbackslash{}( \textbackslash{}vec\{L\}\_i = \textbackslash{}vec\{L\}\_\{1i\} + \textbackslash{}vec\{L\}\_\{2i\} \textbackslash{})) is shown as a vertical arrow pointing upward.
- A circular path is indicated by a blue line, marking the rotation plane.

**After:**

- The same rod with spheres now has an increased length (\textbackslash{}( l\_f = 160 \textbackslash{}, \textbackslash{}text\{cm\} \textbackslash{})).
- No specific angular velocity (\textbackslash{}( \textbackslash{}omega\_f \textbackslash{})) or value is given.
- Linear momenta (\textbackslash{}( \textbackslash{}vec\{p\}\_\{1f\} \textbackslash{}) and \textbackslash{}( \textbackslash{}vec\{p\}\_\{2f\} \textbackslash{})) are depicted as arrows pointing along the new rotation plane.
- The total angular momentum (\textbackslash{}( \textbackslash{}vec\{L\}\_f \textbackslash{})) is again shown with an upward vertical arrow.
- The rotation axis is marked vertically, and the rotation plane is indicated by another blue line, now larger due to the increased radius. 

The scene depicts a conservation of angular momentum scenario with changes in the rotational radius.

\#\# Dataset metadata

Rotational Motion; Physical Model Grounding Reasoning, Multi-Formula Reasoning

\paragraph{Recorded answer/context.}
C: C: 0.20rev/s

\paragraph{Figure/image path.}
/root/proposal\_for\_physic/science-mango/phyx\_mini\_run/phyx\_data/test\_image/711.png

\paragraph{Formalization target.}
create a compiling Lean file with sorry bodies at `PhyXMiniProblems/problem\_phyx\_mini\_0711.lean`. The Lean declarations must preserve the physical quantities, dimensions or dimensional roles, figure labels, governing-law hypotheses, and final relation expressed by this problem.
Use Mathlib/Physlib names found through LeanExplore where available. If a physics API is missing, introduce faithful local abstractions rather than scalar placeholder aliases.

\begin{theorem}[Physics formalization target]
\label{thm:physics:phyx_mini_0711:target}
\lean{PhyXMiniProblems.ProblemPhyXMini0711.finalRotationFrequency_matches_recordedChoiceC}
\uses{def:physics:phyx-mini-0711:phyxminiproblems-problemphyxmini0711-frequencyinrevolutionspersecond, def:physics:phyx-mini-0711:phyxminiproblems-problemphyxmini0711-expandingtwomassrodsetup, def:physics:phyx-mini-0711:phyxminiproblems-problemphyxmini0711-matchesexpandingtwomassrodscenario, def:physics:phyx-mini-0711:phyxminiproblems-problemphyxmini0711-matchesproblemreadouts, def:physics:phyx-mini-0711:phyxminiproblems-problemphyxmini0711-matchessuppliedexpandingrodfigure, def:physics:phyx-mini-0711:phyxminiproblems-problemphyxmini0711-hasphysicalexpandingrodparameters, def:physics:phyx-mini-0711:phyxminiproblems-problemphyxmini0711-satisfiescenteredendpointgeometry, def:physics:phyx-mini-0711:phyxminiproblems-problemphyxmini0711-satisfiestwopointmassrigidrotationlaws, def:physics:phyx-mini-0711:phyxminiproblems-problemphyxmini0711-satisfiesangularmomentumconservationduringexpansion, def:physics:phyx-mini-0711:phyxminiproblems-problemphyxmini0711-recordeddatasetanswer, def:physics:phyx-mini-0711:phyxminiproblems-problemphyxmini0711-matchesdisplayedrotationfrequency, def:physics:phyx-mini-0711:phyxminiproblems-problemphyxmini0711-isuniqueclosestdisplayedrotationfrequency}
The assigned autoformalize agent should translate this physics problem into Lean declarations in the covered file, with theorem and lemma proof bodies written as `by sorry`.
\end{theorem}
\begin{proof}
This is an autoformalization task, not a proof task. Produce statements that can later be proved by the physics prover without weakening the physical model.
\end{proof}

% --- Archon declaration topology begin ---
\section{Declaration topology}
\begin{definition}[rotation Frequency Dimension]
  \label{def:physics:phyx-mini-0711:phyxminiproblems-problemphyxmini0711-rotationfrequencydimension}
  \lean{PhyXMiniProblems.ProblemPhyXMini0711.rotationFrequencyDimension}
  Rotation frequency has inverse-time dimension; one cycle is one revolution
\end{definition}
\begin{proof}
  This is the declared model definition of rotation Frequency Dimension.
\end{proof}
\begin{definition}[moment Of Inertia Dimension]
  \label{def:physics:phyx-mini-0711:phyxminiproblems-problemphyxmini0711-momentofinertiadimension}
  \lean{PhyXMiniProblems.ProblemPhyXMini0711.momentOfInertiaDimension}
  A moment of inertia has dimension mass times length squared
\end{definition}
\begin{proof}
  This is the declared model definition of moment Of Inertia Dimension.
\end{proof}
\begin{definition}[angular Momentum Dimension]
  \label{def:physics:phyx-mini-0711:phyxminiproblems-problemphyxmini0711-angularmomentumdimension}
  \lean{PhyXMiniProblems.ProblemPhyXMini0711.angularMomentumDimension}
  Angular momentum has dimension mass times length squared per time
\end{definition}
\begin{proof}
  This is the declared model definition of angular Momentum Dimension.
\end{proof}
\begin{definition}[Mass Quantity]
  \label{def:physics:phyx-mini-0711:phyxminiproblems-problemphyxmini0711-massquantity}
  \lean{PhyXMiniProblems.ProblemPhyXMini0711.MassQuantity}
  A nonnegative, unit-independent physical mass
\end{definition}
\begin{proof}
  This is the declared model definition of Mass Quantity.
\end{proof}
\begin{definition}[Length Quantity]
  \label{def:physics:phyx-mini-0711:phyxminiproblems-problemphyxmini0711-lengthquantity}
  \lean{PhyXMiniProblems.ProblemPhyXMini0711.LengthQuantity}
  A nonnegative, unit-independent physical length
\end{definition}
\begin{proof}
  This is the declared model definition of Length Quantity.
\end{proof}
\begin{definition}[Rotation Frequency Quantity]
  \label{def:physics:phyx-mini-0711:phyxminiproblems-problemphyxmini0711-rotationfrequencyquantity}
  \lean{PhyXMiniProblems.ProblemPhyXMini0711.RotationFrequencyQuantity}
  \uses{def:physics:phyx-mini-0711:phyxminiproblems-problemphyxmini0711-rotationfrequencydimension}
  A nonnegative ordinary rotation frequency, measured in revolutions per time
\end{definition}
\begin{proof}
  This is the declared model definition of Rotation Frequency Quantity.
\end{proof}
\begin{definition}[Moment Of Inertia Quantity]
  \label{def:physics:phyx-mini-0711:phyxminiproblems-problemphyxmini0711-momentofinertiaquantity}
  \lean{PhyXMiniProblems.ProblemPhyXMini0711.MomentOfInertiaQuantity}
  \uses{def:physics:phyx-mini-0711:phyxminiproblems-problemphyxmini0711-momentofinertiadimension}
  A nonnegative scalar moment of inertia about the displayed rotation axis
\end{definition}
\begin{proof}
  This is the declared model definition of Moment Of Inertia Quantity.
\end{proof}
\begin{definition}[Angular Momentum Magnitude]
  \label{def:physics:phyx-mini-0711:phyxminiproblems-problemphyxmini0711-angularmomentummagnitude}
  \lean{PhyXMiniProblems.ProblemPhyXMini0711.AngularMomentumMagnitude}
  \uses{def:physics:phyx-mini-0711:phyxminiproblems-problemphyxmini0711-angularmomentumdimension}
  A nonnegative angular-momentum magnitude along the displayed upward axis
\end{definition}
\begin{proof}
  This is the declared model definition of Angular Momentum Magnitude.
\end{proof}
\begin{definition}[mass Readout]
  \label{def:physics:phyx-mini-0711:phyxminiproblems-problemphyxmini0711-massreadout}
  \lean{PhyXMiniProblems.ProblemPhyXMini0711.massReadout}
  \uses{def:physics:phyx-mini-0711:phyxminiproblems-problemphyxmini0711-massquantity}
  Read a physical mass in the selected mass unit
\end{definition}
\begin{proof}
  This is the declared model definition of mass Readout.
\end{proof}
\begin{definition}[length Readout]
  \label{def:physics:phyx-mini-0711:phyxminiproblems-problemphyxmini0711-lengthreadout}
  \lean{PhyXMiniProblems.ProblemPhyXMini0711.lengthReadout}
  \uses{def:physics:phyx-mini-0711:phyxminiproblems-problemphyxmini0711-lengthquantity}
  Read a physical length in the selected length unit
\end{definition}
\begin{proof}
  This is the declared model definition of length Readout.
\end{proof}
\begin{definition}[rotation Frequency Readout]
  \label{def:physics:phyx-mini-0711:phyxminiproblems-problemphyxmini0711-rotationfrequencyreadout}
  \lean{PhyXMiniProblems.ProblemPhyXMini0711.rotationFrequencyReadout}
  \uses{def:physics:phyx-mini-0711:phyxminiproblems-problemphyxmini0711-rotationfrequencyquantity}
  Read ordinary rotation frequency in revolutions per selected time unit. The dimensionless cycle counted by this readout is fixed to mean one revolution
\end{definition}
\begin{proof}
  This is the declared model definition of rotation Frequency Readout.
\end{proof}
\begin{definition}[moment Of Inertia Readout]
  \label{def:physics:phyx-mini-0711:phyxminiproblems-problemphyxmini0711-momentofinertiareadout}
  \lean{PhyXMiniProblems.ProblemPhyXMini0711.momentOfInertiaReadout}
  \uses{def:physics:phyx-mini-0711:phyxminiproblems-problemphyxmini0711-momentofinertiaquantity}
  Read a moment of inertia in coherent selected mass and length units
\end{definition}
\begin{proof}
  This is the declared model definition of moment Of Inertia Readout.
\end{proof}
\begin{definition}[angular Momentum Readout]
  \label{def:physics:phyx-mini-0711:phyxminiproblems-problemphyxmini0711-angularmomentumreadout}
  \lean{PhyXMiniProblems.ProblemPhyXMini0711.angularMomentumReadout}
  \uses{def:physics:phyx-mini-0711:phyxminiproblems-problemphyxmini0711-angularmomentummagnitude}
  Read angular momentum in coherent selected mechanical units
\end{definition}
\begin{proof}
  This is the declared model definition of angular Momentum Readout.
\end{proof}
\begin{definition}[length In Centimeters]
  \label{def:physics:phyx-mini-0711:phyxminiproblems-problemphyxmini0711-lengthincentimeters}
  \lean{PhyXMiniProblems.ProblemPhyXMini0711.lengthInCentimeters}
  \uses{def:physics:phyx-mini-0711:phyxminiproblems-problemphyxmini0711-lengthquantity, def:physics:phyx-mini-0711:phyxminiproblems-problemphyxmini0711-lengthreadout}
  Centimetre readout of a physical length
\end{definition}
\begin{proof}
  This is the declared model definition of length In Centimeters.
\end{proof}
\begin{definition}[frequency In Revolutions Per Second]
  \label{def:physics:phyx-mini-0711:phyxminiproblems-problemphyxmini0711-frequencyinrevolutionspersecond}
  \lean{PhyXMiniProblems.ProblemPhyXMini0711.frequencyInRevolutionsPerSecond}
  \uses{def:physics:phyx-mini-0711:phyxminiproblems-problemphyxmini0711-rotationfrequencyquantity, def:physics:phyx-mini-0711:phyxminiproblems-problemphyxmini0711-rotationfrequencyreadout}
  Revolutions-per-second readout of an ordinary rotation frequency
\end{definition}
\begin{proof}
  This is the declared model definition of frequency In Revolutions Per Second.
\end{proof}
\begin{definition}[Configuration]
  \label{def:physics:phyx-mini-0711:phyxminiproblems-problemphyxmini0711-configuration}
  \lean{PhyXMiniProblems.ProblemPhyXMini0711.Configuration}
  The two configurations shown in the panels labelled Before and After
\end{definition}
\begin{proof}
  This is the declared model definition of Configuration.
\end{proof}
\begin{definition}[Endpoint Mass]
  \label{def:physics:phyx-mini-0711:phyxminiproblems-problemphyxmini0711-endpointmass}
  \lean{PhyXMiniProblems.ProblemPhyXMini0711.EndpointMass}
  The two material spheres labelled 1 and 2 in both panels
\end{definition}
\begin{proof}
  This is the declared model definition of Endpoint Mass.
\end{proof}
\begin{definition}[Rod Mass Model]
  \label{def:physics:phyx-mini-0711:phyxminiproblems-problemphyxmini0711-rodmassmodel}
  \lean{PhyXMiniProblems.ProblemPhyXMini0711.RodMassModel}
  The idealized contribution of the connecting rod to the mass distribution
\end{definition}
\begin{proof}
  This is the declared model definition of Rod Mass Model.
\end{proof}
\begin{definition}[Expansion Mechanism]
  \label{def:physics:phyx-mini-0711:phyxminiproblems-problemphyxmini0711-expansionmechanism}
  \lean{PhyXMiniProblems.ProblemPhyXMini0711.ExpansionMechanism}
  The mechanism stated to change the rod length
\end{definition}
\begin{proof}
  This is the declared model definition of Expansion Mechanism.
\end{proof}
\begin{definition}[Rotation Axis Geometry]
  \label{def:physics:phyx-mini-0711:phyxminiproblems-problemphyxmini0711-rotationaxisgeometry}
  \lean{PhyXMiniProblems.ProblemPhyXMini0711.RotationAxisGeometry}
  Relation of the fixed rotation axis to the straight rod
\end{definition}
\begin{proof}
  This is the declared model definition of Rotation Axis Geometry.
\end{proof}
\begin{definition}[Figure Label]
  \label{def:physics:phyx-mini-0711:phyxminiproblems-problemphyxmini0711-figurelabel}
  \lean{PhyXMiniProblems.ProblemPhyXMini0711.FigureLabel}
  Literal labels visible in image 711.png
\end{definition}
\begin{proof}
  This is the declared model definition of Figure Label.
\end{proof}
\begin{definition}[Expanding Rod Figure]
  \label{def:physics:phyx-mini-0711:phyxminiproblems-problemphyxmini0711-expandingrodfigure}
  \lean{PhyXMiniProblems.ProblemPhyXMini0711.ExpandingRodFigure}
  \uses{def:physics:phyx-mini-0711:phyxminiproblems-problemphyxmini0711-configuration, def:physics:phyx-mini-0711:phyxminiproblems-problemphyxmini0711-endpointmass, def:physics:phyx-mini-0711:phyxminiproblems-problemphyxmini0711-figurelabel}
  Literal and qualitative information transcribed from the supplied raster. Pixel radii describe only the perspective drawing; they are not physical length measurements. The final omega label contains no numerical answer
\end{definition}
\begin{proof}
  This is the declared model definition of Expanding Rod Figure.
\end{proof}
\begin{definition}[Expanding Two Mass Rod Setup]
  \label{def:physics:phyx-mini-0711:phyxminiproblems-problemphyxmini0711-expandingtwomassrodsetup}
  \lean{PhyXMiniProblems.ProblemPhyXMini0711.ExpandingTwoMassRodSetup}
  \uses{def:physics:phyx-mini-0711:phyxminiproblems-problemphyxmini0711-massquantity, def:physics:phyx-mini-0711:phyxminiproblems-problemphyxmini0711-lengthquantity, def:physics:phyx-mini-0711:phyxminiproblems-problemphyxmini0711-rotationfrequencyquantity, def:physics:phyx-mini-0711:phyxminiproblems-problemphyxmini0711-momentofinertiaquantity, def:physics:phyx-mini-0711:phyxminiproblems-problemphyxmini0711-angularmomentummagnitude, def:physics:phyx-mini-0711:phyxminiproblems-problemphyxmini0711-configuration, def:physics:phyx-mini-0711:phyxminiproblems-problemphyxmini0711-endpointmass, def:physics:phyx-mini-0711:phyxminiproblems-problemphyxmini0711-rodmassmodel, def:physics:phyx-mini-0711:phyxminiproblems-problemphyxmini0711-expansionmechanism, def:physics:phyx-mini-0711:phyxminiproblems-problemphyxmini0711-rotationaxisgeometry, def:physics:phyx-mini-0711:phyxminiproblems-problemphyxmini0711-expandingrodfigure}
  The final rotation frequency, both moments of inertia, and both total angular momenta are independent observables. Subsequent predicates relate them by general geometry and mechanics laws; no field contains the requested result
\end{definition}
\begin{proof}
  This is the declared model definition of Expanding Two Mass Rod Setup.
\end{proof}
\begin{definition}[Matches Expanding Two Mass Rod Scenario]
  \label{def:physics:phyx-mini-0711:phyxminiproblems-problemphyxmini0711-matchesexpandingtwomassrodscenario}
  \lean{PhyXMiniProblems.ProblemPhyXMini0711.MatchesExpandingTwoMassRodScenario}
  \uses{def:physics:phyx-mini-0711:phyxminiproblems-problemphyxmini0711-expandingtwomassrodsetup}
  The prose-level two-equal-mass, massless, gas-expanding rod scenario
\end{definition}
\begin{proof}
  This is the declared model definition of Matches Expanding Two Mass Rod Scenario.
\end{proof}
\begin{definition}[Matches Problem Readouts]
  \label{def:physics:phyx-mini-0711:phyxminiproblems-problemphyxmini0711-matchesproblemreadouts}
  \lean{PhyXMiniProblems.ProblemPhyXMini0711.MatchesProblemReadouts}
  \uses{def:physics:phyx-mini-0711:phyxminiproblems-problemphyxmini0711-lengthincentimeters, def:physics:phyx-mini-0711:phyxminiproblems-problemphyxmini0711-frequencyinrevolutionspersecond, def:physics:phyx-mini-0711:phyxminiproblems-problemphyxmini0711-expandingtwomassrodsetup}
  Numerical physical readouts stated in the problem. There is deliberately no field for the final rotation frequency
\end{definition}
\begin{proof}
  This is the declared model definition of Matches Problem Readouts.
\end{proof}
\begin{definition}[Matches Supplied Expanding Rod Figure]
  \label{def:physics:phyx-mini-0711:phyxminiproblems-problemphyxmini0711-matchessuppliedexpandingrodfigure}
  \lean{PhyXMiniProblems.ProblemPhyXMini0711.MatchesSuppliedExpandingRodFigure}
  \uses{def:physics:phyx-mini-0711:phyxminiproblems-problemphyxmini0711-expandingtwomassrodsetup}
  Evidence read directly from image 711.png: labels, rods, spheres, circular paths, tangential linear-momentum arrows, upward total-angular-momentum arrows, and the printed before/after data. The post-expansion frequency is shown only as the unknown symbol omega-f
\end{definition}
\begin{proof}
  This is the declared model definition of Matches Supplied Expanding Rod Figure.
\end{proof}
\begin{definition}[Has Physical Expanding Rod Parameters]
  \label{def:physics:phyx-mini-0711:phyxminiproblems-problemphyxmini0711-hasphysicalexpandingrodparameters}
  \lean{PhyXMiniProblems.ProblemPhyXMini0711.HasPhysicalExpandingRodParameters}
  \uses{def:physics:phyx-mini-0711:phyxminiproblems-problemphyxmini0711-massreadout, def:physics:phyx-mini-0711:phyxminiproblems-problemphyxmini0711-lengthreadout, def:physics:phyx-mini-0711:phyxminiproblems-problemphyxmini0711-rotationfrequencyreadout, def:physics:phyx-mini-0711:phyxminiproblems-problemphyxmini0711-momentofinertiareadout, def:physics:phyx-mini-0711:phyxminiproblems-problemphyxmini0711-angularmomentumreadout, def:physics:phyx-mini-0711:phyxminiproblems-problemphyxmini0711-expandingtwomassrodsetup}
  Positivity and nondegeneracy conditions for the rotating physical branch
\end{definition}
\begin{proof}
  This is the declared model definition of Has Physical Expanding Rod Parameters.
\end{proof}
\begin{definition}[Satisfies Centered Endpoint Geometry]
  \label{def:physics:phyx-mini-0711:phyxminiproblems-problemphyxmini0711-satisfiescenteredendpointgeometry}
  \lean{PhyXMiniProblems.ProblemPhyXMini0711.SatisfiesCenteredEndpointGeometry}
  \uses{def:physics:phyx-mini-0711:phyxminiproblems-problemphyxmini0711-lengthreadout, def:physics:phyx-mini-0711:phyxminiproblems-problemphyxmini0711-expandingtwomassrodsetup}
  Each endpoint mass lies half a rod length from the midpoint axis
\end{definition}
\begin{proof}
  This is the declared model definition of Satisfies Centered Endpoint Geometry.
\end{proof}
\begin{definition}[Satisfies Two Point Mass Rigid Rotation Laws]
  \label{def:physics:phyx-mini-0711:phyxminiproblems-problemphyxmini0711-satisfiestwopointmassrigidrotationlaws}
  \lean{PhyXMiniProblems.ProblemPhyXMini0711.SatisfiesTwoPointMassRigidRotationLaws}
  \uses{def:physics:phyx-mini-0711:phyxminiproblems-problemphyxmini0711-massreadout, def:physics:phyx-mini-0711:phyxminiproblems-problemphyxmini0711-lengthreadout, def:physics:phyx-mini-0711:phyxminiproblems-problemphyxmini0711-rotationfrequencyreadout, def:physics:phyx-mini-0711:phyxminiproblems-problemphyxmini0711-momentofinertiareadout, def:physics:phyx-mini-0711:phyxminiproblems-problemphyxmini0711-angularmomentumreadout, def:physics:phyx-mini-0711:phyxminiproblems-problemphyxmini0711-endpointmass, def:physics:phyx-mini-0711:phyxminiproblems-problemphyxmini0711-expandingtwomassrodsetup}
  General rigid-rotation laws for two point masses on a massless rod. The factor 2 * pi converts ordinary rotation frequency in revolutions per unit time to angular speed in radians per unit time. These laws are uniform in configuration and contain no specialized final-frequency value
\end{definition}
\begin{proof}
  This is the declared model definition of Satisfies Two Point Mass Rigid Rotation Laws.
\end{proof}
\begin{definition}[Satisfies Angular Momentum Conservation During Expansion]
  \label{def:physics:phyx-mini-0711:phyxminiproblems-problemphyxmini0711-satisfiesangularmomentumconservationduringexpansion}
  \lean{PhyXMiniProblems.ProblemPhyXMini0711.SatisfiesAngularMomentumConservationDuringExpansion}
  \uses{def:physics:phyx-mini-0711:phyxminiproblems-problemphyxmini0711-angularmomentumreadout, def:physics:phyx-mini-0711:phyxminiproblems-problemphyxmini0711-expandingtwomassrodsetup}
  With zero external torque about the fixed midpoint axis, total angular momentum is the same before and after the gas-driven expansion. This premise does not specify the final frequency
\end{definition}
\begin{proof}
  This is the declared model definition of Satisfies Angular Momentum Conservation During Expansion.
\end{proof}
\begin{lemma}[midpoint Moment Of Inertia readout]
  \label{lem:physics:phyx-mini-0711:phyxminiproblems-problemphyxmini0711-midpointmomentofinertia-readout}
  \lean{PhyXMiniProblems.ProblemPhyXMini0711.midpointMomentOfInertia_readout}
  \uses{def:physics:phyx-mini-0711:phyxminiproblems-problemphyxmini0711-massreadout, def:physics:phyx-mini-0711:phyxminiproblems-problemphyxmini0711-lengthreadout, def:physics:phyx-mini-0711:phyxminiproblems-problemphyxmini0711-momentofinertiareadout, def:physics:phyx-mini-0711:phyxminiproblems-problemphyxmini0711-expandingtwomassrodsetup, def:physics:phyx-mini-0711:phyxminiproblems-problemphyxmini0711-matchesexpandingtwomassrodscenario, def:physics:phyx-mini-0711:phyxminiproblems-problemphyxmini0711-satisfiescenteredendpointgeometry, def:physics:phyx-mini-0711:phyxminiproblems-problemphyxmini0711-satisfiestwopointmassrigidrotationlaws}
  Two equal endpoint masses at radius l/2 have I = m l\textasciicircum{}2 / 2
\end{lemma}
\begin{proof}
  The conclusion follows from the governing relations and figure readouts named in the statement of midpoint Moment Of Inertia readout.
\end{proof}
\begin{lemma}[rod Length Squared mul frequency conserved]
  \label{lem:physics:phyx-mini-0711:phyxminiproblems-problemphyxmini0711-rodlengthsquared-mul-frequency-conserved}
  \lean{PhyXMiniProblems.ProblemPhyXMini0711.rodLengthSquared_mul_frequency_conserved}
  \uses{def:physics:phyx-mini-0711:phyxminiproblems-problemphyxmini0711-lengthincentimeters, def:physics:phyx-mini-0711:phyxminiproblems-problemphyxmini0711-frequencyinrevolutionspersecond, def:physics:phyx-mini-0711:phyxminiproblems-problemphyxmini0711-expandingtwomassrodsetup, def:physics:phyx-mini-0711:phyxminiproblems-problemphyxmini0711-matchesexpandingtwomassrodscenario, def:physics:phyx-mini-0711:phyxminiproblems-problemphyxmini0711-hasphysicalexpandingrodparameters, def:physics:phyx-mini-0711:phyxminiproblems-problemphyxmini0711-satisfiescenteredendpointgeometry, def:physics:phyx-mini-0711:phyxminiproblems-problemphyxmini0711-satisfiestwopointmassrigidrotationlaws, def:physics:phyx-mini-0711:phyxminiproblems-problemphyxmini0711-satisfiesangularmomentumconservationduringexpansion}
  Conservation of angular momentum cancels the common mass and 2 pi factors, leaving l\_i\textasciicircum{}2 f\_i = l\_f\textasciicircum{}2 f\_f in centimetre/second readouts
\end{lemma}
\begin{proof}
  The conclusion follows from the governing relations and figure readouts named in the statement of rod Length Squared mul frequency conserved.
\end{proof}
\begin{lemma}[final Rotation Frequency exact]
  \label{lem:physics:phyx-mini-0711:phyxminiproblems-problemphyxmini0711-finalrotationfrequency-exact}
  \lean{PhyXMiniProblems.ProblemPhyXMini0711.finalRotationFrequency_exact}
  \uses{def:physics:phyx-mini-0711:phyxminiproblems-problemphyxmini0711-frequencyinrevolutionspersecond, def:physics:phyx-mini-0711:phyxminiproblems-problemphyxmini0711-expandingtwomassrodsetup, def:physics:phyx-mini-0711:phyxminiproblems-problemphyxmini0711-matchesexpandingtwomassrodscenario, def:physics:phyx-mini-0711:phyxminiproblems-problemphyxmini0711-matchesproblemreadouts, def:physics:phyx-mini-0711:phyxminiproblems-problemphyxmini0711-hasphysicalexpandingrodparameters, def:physics:phyx-mini-0711:phyxminiproblems-problemphyxmini0711-satisfiescenteredendpointgeometry, def:physics:phyx-mini-0711:phyxminiproblems-problemphyxmini0711-satisfiestwopointmassrigidrotationlaws, def:physics:phyx-mini-0711:phyxminiproblems-problemphyxmini0711-satisfiesangularmomentumconservationduringexpansion}
  Using the stated lengths and initial frequency gives the exact physical readout 2 * (50 / 160)\textasciicircum{}2 = 25 / 128 rev/s
\end{lemma}
\begin{proof}
  The conclusion follows from the governing relations and figure readouts named in the statement of final Rotation Frequency exact.
\end{proof}
\begin{definition}[Answer Choice]
  \label{def:physics:phyx-mini-0711:phyxminiproblems-problemphyxmini0711-answerchoice}
  \lean{PhyXMiniProblems.ProblemPhyXMini0711.AnswerChoice}
  Labels of the four rotation-frequency choices in the source
\end{definition}
\begin{proof}
  This is the declared model definition of Answer Choice.
\end{proof}
\begin{definition}[revolutions Per Second]
  \label{def:physics:phyx-mini-0711:phyxminiproblems-problemphyxmini0711-answerchoice-revolutionspersecond}
  \lean{PhyXMiniProblems.ProblemPhyXMini0711.AnswerChoice.revolutionsPerSecond}
  \uses{def:physics:phyx-mini-0711:phyxminiproblems-problemphyxmini0711-answerchoice}
  Revolutions-per-second value printed beside each answer label
\end{definition}
\begin{proof}
  This is the declared model definition of revolutions Per Second.
\end{proof}
\begin{definition}[recorded Dataset Answer]
  \label{def:physics:phyx-mini-0711:phyxminiproblems-problemphyxmini0711-recordeddatasetanswer}
  \lean{PhyXMiniProblems.ProblemPhyXMini0711.recordedDatasetAnswer}
  \uses{def:physics:phyx-mini-0711:phyxminiproblems-problemphyxmini0711-answerchoice}
  The answer label recorded by the supplied dataset metadata
\end{definition}
\begin{proof}
  This is the declared model definition of recorded Dataset Answer.
\end{proof}
\begin{definition}[Matches Displayed Rotation Frequency]
  \label{def:physics:phyx-mini-0711:phyxminiproblems-problemphyxmini0711-matchesdisplayedrotationfrequency}
  \lean{PhyXMiniProblems.ProblemPhyXMini0711.MatchesDisplayedRotationFrequency}
  \uses{def:physics:phyx-mini-0711:phyxminiproblems-problemphyxmini0711-frequencyinrevolutionspersecond, def:physics:phyx-mini-0711:phyxminiproblems-problemphyxmini0711-expandingtwomassrodsetup, def:physics:phyx-mini-0711:phyxminiproblems-problemphyxmini0711-answerchoice}
  Agreement with a value printed to two decimal places: the exact physical readout differs from the displayed value by less than half of 0.01 rev/s
\end{definition}
\begin{proof}
  This is the declared model definition of Matches Displayed Rotation Frequency.
\end{proof}
\begin{definition}[Is Unique Closest Displayed Rotation Frequency]
  \label{def:physics:phyx-mini-0711:phyxminiproblems-problemphyxmini0711-isuniqueclosestdisplayedrotationfrequency}
  \lean{PhyXMiniProblems.ProblemPhyXMini0711.IsUniqueClosestDisplayedRotationFrequency}
  \uses{def:physics:phyx-mini-0711:phyxminiproblems-problemphyxmini0711-frequencyinrevolutionspersecond, def:physics:phyx-mini-0711:phyxminiproblems-problemphyxmini0711-expandingtwomassrodsetup, def:physics:phyx-mini-0711:phyxminiproblems-problemphyxmini0711-answerchoice}
  A displayed choice is strictly closer than every other displayed value
\end{definition}
\begin{proof}
  This is the declared model definition of Is Unique Closest Displayed Rotation Frequency.
\end{proof}
% --- Archon declaration topology end ---
% --- Archon physics formalization source end ---
